\cotefile{156-4}
\scanpage{26}\mspage{25}
On reprend tout !

\uline{Définition-\unsure{scholie}} Une \uline{contrée} \inserted{$C$} est \unsure{definie|définie} par un ensemble $L$ (l'ens. des \uline{lieux} de la contrée), et la donnée de certains ensembles de parties de $L$, appelés les \uline{figures} \inserted{(de la contrée)}.

L'ens. $\unsure{\mathcal{F}|\mathscr{F}}$ des figures de $C$ est donc une \deleted{ensemble} partie de $P(P(L))$ (donc $F\in P(P(P(L)))$). Si \deleted{$\unsure{\mathcal{F}|F}$}$F$ est une figure, ses éléments\textsuperscript{\unsure{$\times$/}} (qui sont des parties de $L$) s'appellent les \uline{strates} (\uline{fermées}) de la figure $F$. \unsure{Là|Lá|La} réunion des strates \deleted{\unsure{fermées|ferm}} de $F$ est appelée le \uline{support} de $F$, \inserted{noté $|F|$.} On pose des \uline{axiomes des contrées}, qui vont être dégagés par la suite.

\uline{Cont}~1 \quad a) \unsure{La|la} partie vide de $P(L)$ est une \deleted{strate} figure, appelée ``\uline{figure vide}''. C'est donc la figure unique qui n'a pas de strates.

% slanted
\margin{\unsure{Contenu|Contenue} dans \uline{Cont} \unsure{4|5} plus bas}
\margin{$\big\Vert$}
b) $\forall\,x\in L$, \deleted{\unsure{la|le}}\inserted{il y a une} figure ayant \deleted{\unsure{la|les}} strate unique $\{x\}$. On l'appelle \uline{figure ponctuelle} \unsure{definie|définie} par $x$
